\chapter{Appendix 7: Smart Card Technology}

\section{Introduction}

Smart cards represent a new technology that has tremendous potential
for enhancing the security of distributed systems. This appendix
introduces smart card terminology, describes current smart card
technology (including physical and logical characteristics), discusses
the current status of smart card standardization, and closes with a
brief discussion of the benefits a company could gain by utilizing
smart cards.


\section{Terminology}

\begin{htmllist}
  \htmlitemmark{RedBall}

 \item[Account]
 
 A security database entry associating a principal with a key and other
 security-related information such as privileges.

 \item[Issuer]
 
 In this RFC series, an organization that purchases smart cards and issues
 them to its users for use in accessing the issuer's services or systems.

 \item[Dedicated file]
 
 A smart card file containing file control information and, optionally,
 memory available for allocation. It may be the parent of elementary files
 and/or dedicated files ISO 7816-4; roughly, a smart card ``file
 directory''.

 \item[EEPROM]

 Electrically Erasable Programmable Read-Only Memory.

 \item[Elementary file]
 
 A set of data units or records in the smart card that share the same
 identifier ISO 7816-4; roughly, a smart card ``file''.

 \item[File control information (FCI)]
 
 Logical, structural and security attributes of a file; a string of data
 bytes available in response to a SELECT FILE command ISO 7816-4.

 \item[FRAM]

 Ferro-Electric Random Access Memory.

 \item[GSM]
 
 Groupe Special Mobile; European digital cellular telephone standard,
 providing security services such as user authentication, traffic
 confidentiality, and key distribution.

 \item[IC]
 
 Integrated Circuit. A single circuit containing multiple transistors.

 \item[Interactive principal]

 A human user of a system. See ``principal''.

 \item[Key]
   
 A parameter used in conjunction with a cryptographic algorithm that
 determines:

 \begin{itemize}
   \item The transformation of plaintext data into ciphertext data,
   \item The transformation of ciphertext data into plaintext data,
   \item A digital signature, or
   \item A message authentication code.
 \end{itemize}

 \item[Long-term key]
  
 A key that is used over a relatively long period of time. Compare with
 ``session key''. Each principal can have a unique long-term key that is
 used in authenticating a user or server claiming to be that principal.
 For interactive principals, the long-term key is derived from the user's
 password.

 \item[Master file]
  
 The mandatory unique dedicated file representing the root of the smart
 card file structure ISO 7816-4.

 \item[Message authentication code (MAC)]
  
 A cryptographic checksum, based on DES.

 \item[Noninteractive principal]
  
 A principal that is an instance of a company server, an instance of a
 company application server, a computer in a company cell, or a cell (more
 correctly, ``an Authentication Service surrogate''). See ``principal''.

 \item[Passivation layer]
  
 A layer of dielectric material covering the circuitry on a smart card
 that protects the chip from impurities and dust and prevents passage of
 radiation associated with probes, such as electron-beam microscopy. Often
 referred to as a ``seal''.

 \item[Password]
  
 A sequence of alphanumeric and punctuation characters entered by a user
 to authenticate to a computer system (including to a smart card).

 \item[PIN]
  
 Personal Identification Number; a 4 to 12 character alphanumeric code or
 password used to authenticate a person's identity, commonly used in
 banking applications. In connection with smart cards, this document uses
 the more general term ``password'' rather than ``PIN''.

 \item[Principal]
  
 An entity that is capable of believing that it can communicate securely
 with another entity. In a company Security, principals are represented as
 entries in the Registry database. See ``interactive principal'' and
 ``noninteractive principal''.

 \item[Session key]
  
 A random number generated to serve as a key for a specific transaction or
 set of transactions.
  
 \item[Smart card]
  
 A credit card-sized, tamper-resistant security device that relies on VLSI
 chip technology for information storage and information processing.

 \item[User]
  
 A person who attempts to access a computer system; a person who has been
 issued a smart card.

 \item[VLSI]
  
 Very Large Scale Integration; integration of thousands or more transitors
 on a single chip, enabling single-chip implementations of CPU, RAM, ROM,
 etc.
\end{htmllist}


\section{Smart Card Technology}

\subsection{Basic Concepts}

A smart card is a credit card-sized, tamper-resistant security device
that offers functions for secure information storage and information
processing that rely on VLSI chip technology. A smart card actually
contains a secure microprocessor chip embedded in the card. The chip
can implement a secure file system, compute cryptographic functions,
and actively detect invalid access attempts. With proper application of
file system access rights, a smart card can be safely used by multiple,
independent applications.\\

A smart card is distinguished from a magnetic stripe card (e.g.,
typical U.S. credit card) in that a magnetic stripe card has no VLSI
circuitry, and thus no active security procedures and no built-in
tamper-resistance. Anyone with an appropriate card reader can read
whatever is on the card. The current cost for magnetic stripe cards
(March 1994) is around \$.50-\$1.00 per card in quantities of 1000.\\

Smart cards are also distinguished from ``SuperSmart'' or ``token''
cards. These use a small LCD to display the time of day encrypted using
an internal key that is programmed at the time of manufacture. Such
cards have no storage capabilities and no interface to external
computer systems. They perform a single, well-defined function that can
be used in authentication protocols: users manually type the digits
displayed on the card (which are updated periodically) as part of
authenticating to a host. These cards are not capable of performing
more extensive functions as part of a security system, and thus are not
evolutionary. Cost (as of March 1994) for ``SuperSmart'' cards is on
the order of \$60/card in quantities of 250-500 (based on a cost of \$20
per year, and battery lifetime of around 3 years).

\subsection{Physical Characteristics}

\begin{htmllist}
  \htmlitemmark{RedBall}

  \latexhtml{\item {\bf Card types}}{\item[Card types]}
  
  Within the family of smart cards, there are two general categories: cards
  with contacts and cards without contacts. Contact cards contain physical
  contact points on the surface of the card that allow transmission of
  commands, data, and status information between the card and a card
  reader. Contactless cards also require physical contact with a reader,
  but get their power via induction rather than through one of the
  contacts. This difference affects only the very lowest level (physical
  level) of the protocols used with smart cards.


  \latexhtml{\item {\bf Physical specifications}}{\item[Physical
  specifications]}
  
  Focusing now on contact smart cards, potential uses depend on card
  specifications. Microprocessor and memory technology have steep
  technology curves, but the following specifications of typical available
  contact smart cards can serve as a current baseline (March 1994):

  \latexhtml{\begin{itemize}}{\begin{htmllist}
        \htmlitemmark{WhiteBall}}

      \latexhtml{\item Clock rate:}{\item[Clock rate]} 3.5-8 MHZ.

      \latexhtml{\item Non-volatile memory:}{\item[Non-volatile memory]} 8-16
      Kbytes (for data storage).

      \latexhtml{\item ROM:}{\item[ROM]} 8-16 Kbytes (for the card operating
      system).

      \latexhtml{\item RAM:}{\item[RAM]} Around 256-500 bytes of RAM (for
      operating systemcomputation).

      \latexhtml{\item EEPROM:}{\item[EEPROM]} 2-8 Kbytes
      (externally-accessible, non-volatile).

      \latexhtml{\item Write cycles:}{\item[Write cycles]} Approximately
      104-105.
      
      \latexhtml{\item Processing time:}{\item[Processing time]} Average of
      about 60 milliseconds (including transmission time) to retrieve a
      56-bit number from a protected elementary file on the card. Depends
      primarily on transmission time between the host system and the card
      reader, so varies depending on capacity of that line.
      
      \latexhtml{\item Cost:}{\item[Cost]} On the order of \$7-\$12/card in
      quantities of 1000-5000; card reader cost is on the order of
      \$150-200 each in quantities of 100-500.
  \latexhtml{\end{itemize}}{\end{htmllist}}

  \htmlitemmark{RedBall}
  
  In the next few years, 16-32 bit RISC processors running at 20 MHz are
  likely to become available on smart cards. This will define a new
  generation of smart cards, as the power requirements and frequency
  radiation of these processors can not be handled within the current ISO
  7816-1,2 standard format. By 1995, it is expected that FRAM technology
  will be available, supporting more write cycles (108). This technology
  may also be incorporated into the new generation, although currently it
  is too expensive.

  \latexhtml{\item {\bf Physical security characteristics}}{\item[Physical
    security characteristics]}
  
  Smart cards incorporate physical tamper-resistance circuitry that
  responds to tampering by inhibiting the output function. There is a
  dielectric ``passivation layer'' covering the chip. The passivation layer
  protects the chip from impurities and dust, and prevents passage of
  radiation associated with probes, including electron-beam microscopy. The
  circuitry is capable of reacting to light (indicating the passivation
  layer has been broken); temperature, voltage, and frequency fluctuations
  outside the specified operating range. There are physical memory
  protection mechanisms, including memory scrambling, which make reverse
  engineering more difficult and hinders an attack trying to erase a
  selected data item in memory. Fuses are used during the manufacturing
  cycle to permanently disable ``test'' mode(s) once tests have been passed
  and the card is ready for distribution to issuers.\\

  Smart cards can be manufactured with varying levels of physical security
  features, with higher costs associated with higher levels of security.
  Purchasers must evaluate with the vendor the level of physical security
  appropriate to their application.
\end{htmllist}

\subsection{Logical Characteristics}

The smart card processing unit executes a chip operating system that
implements a hierarchical file system on the non-volatile memory of the
card, and a set of access and control operations both for the card
itself and for the file system.\\

The following descriptions of card data structures, card operations,
and card security architecture are taken from ISO 7816-4. This
document, although still a draft, is currently undergoing balloting
through ISO. It is complete enough to indicate the general types of
structures and operations that cards should be expected to support.

\begin{htmllist}
  \htmlitemmark{RedBall}

    \latexhtml{\item {\bf Data structures}}{\item[Data structures]}
    
    The file system supports a special root directory file (``master
    file''), optional sub-directory files (``dedicated files''), and data
    files (``elementary files''). The identifiers of all files along the
    path from the master file down to a specific file unambiguously
    identify the specific file. All three categories of files contain
    control information, such as a file identifier, file name,
    specifications of record or data lengths in the file, etc. The draft
    standards specify various types of elementary file structures: a
    sequence of records of identical length, a sequence of records of
    variable length, a sequence of records with identical length organized
    as a ring, and a ``transparent structure'' that is seen at the
    interface as a sequence of data units.


    \latexhtml{\item {\bf Operations}}{\item[Operations]}
    
    The following commands are a subset of those specified in ISO 7816-4,
    and are included here merely to give a basic understanding of the types
    of operations supported on smart cards.

  \latexhtml{\begin{itemize}}{\begin{htmllist}
        \htmlitemmark{WhiteBall}}
      
    \latexhtml{\item READ BINARY:}{\item[READ BINARY]} Return all or part
        of the contents of an elementary file.
      
    \latexhtml{\item WRITE BINARY:}{\item[WRITE BINARY]} Write binary
    values into an elementary file.
    
    \latexhtml{\item ERASE BINARY:}{\item[ERASE BINARY]} Erase all or part of
      the contents an elementary file.
      
    \latexhtml{\item SELECT FILE:}{\item[SELECT FILE]} Set a ``current file''
      that may be referred to implicitly in subsequent commands.
      
    \latexhtml{\item VERIFY:}{\item[VERIFY]} Compare verification data sent
      from the host system with reference data stored in the card (e.g., a
      password).
      
    \latexhtml{\item INTERNAL AUTHENTICATE:}{\item[INTERNAL AUTHENTICATE]}
      Compute authentication data using challenge data sent from the host
      system and a relevant secret (e.g., a key) stored in the card.
      
    \latexhtml{\item EXTERNAL AUTHENTICATE:}{\item[EXTERNAL AUTHENTICATE]} 
      Conditionally update the security status of the card based on the
      result of a computation based on a challenge previously issued by the
      card, a key stored in the card, and authentication data sent from the
      host system.
  \latexhtml{\end{itemize}}{\end{htmllist}}

  \htmlitemmark{RedBall}

  \latexhtml{\item {\bf Security Architecture}}{\item[Security Architecture]}
  
  Smart cards implement three levels of logical access control. The first
  is the association of a set of privileges with a user's password, and the
  ability to control access to files on the card based on those privileges.
  The second level of logical access control is the ability to detect and
  respond to a sequence of invalid access attempts. The third level is the
  ``logical channel'' -- a logical link between the host system and a file
  on the smart card.

  \latexhtml{\begin{itemize}}{\begin{htmllist}
        \htmlitemmark{WhiteBall}}

  \latexhtml{\item Privileges and access control}{\item[Privileges and
    access control]}
  
  Two categories of access control mechanisms are promoted today in the
  smart card market. Both mechanisms are built-in characteristics of the
  relation between privileges and users for given objects.

  \begin{itemize}
  \item One system is the combinational system, where privileges are a
    first order logic function between authorization variables (v[i]). For
    example f(v[1], v[2], ..., v[n], p1) will be TRUE for a user/subject
    for a READ access, if and only if the user/subject knows the values of
    every authorization variable v[i] AND knows the password p1. The
    function f is combinational, e.g., (v[1] AND v[2] OR v[3] XOR v[4]) AND
    value\_in\_pwd\_filex\_of(p1).
    
  \item The second system is a sequential system where privileges are the
    result of a sequential function between authorization variables (v[i]).
    For example, f(v[1], v[2], ..., v[n], p1) gives READ access for an
    object to a user/subject, if and only if the user/subject knows which
    values (v[i]) have to be presented in state s[k], where k identifies a
    sequential file selection and key presentation step.s[k] is determined
    by the current selected file and the current keys activated.
  \end{itemize}

  \latexhtml{\item Invalid access attempts}{\item[Invalid access attempts]}
  
  Most, but not all, smart cards keep a record of sequential invalid access
  attempts (a supplied set of parameters for a function that fail to
  evaluate to TRUE), and deny further access to the card (or to the
  targetted file) once the count reaches a certain limit. In some cards the
  limit is configurable, while in others it is fixed at a small number such
  as 3 or 7. The count is reset to 0 when a valid access is made. Exceeding
  the limit either invalidates the card entirely, or puts it in a state
  where only a limited set of operations is available.  These limited
  operations may be sufficient for an administrator to restore access.\\
  
  Denial of access after a small number of invalid attempts prevents
  ``password guessing'' attacks on the card.


  \latexhtml{\item Logical channels}{\item[Logical channels]}
  
  A ``logical channel'' is a logical link between the host system and a
  file on the smart card, either the Master File, a Dedicated File, or an
  Elementary File. When logical channels are in use, the selection of a
  file associates the file and its security status with the logical channel
  encoded in a reserved field of the selection command header.\\

  Logical channels provide a mechanism for allowing multiple, independent
  applications to use the storage capabilities of the card. The card
  interface software on the host system must manage the mapping between
  processes and logical channels; the channel numbers are either assigned
  by the external world or by the card itself.
  \latexhtml{\end{itemize}}{\end{htmllist}}


The logical channel portion of ISO 7816-4 conveys 2 concepts. The first one
deals with a logical link to files and requires the outside world to manage
the channel numbers. In the second concept, the card allocates the logical
channel number and supports a mechanism similar to swapping with a stack
number. In both situations, cards implementing these standards grant
applications the same control of access to files and data structures
(without losing the security status) as if only one application had access.
\end{htmllist}

\subsection{Life Cycle}

As stand-alone active entities, smart cards have a life-cycle and
actors.\\

First, some terminology. The ``chip manufacturer'' is responsible for
the chip design, production, and testing. The ``card manufacturer''
embeds the chip into the plastic card. The ``card issuer'' creates,
designs, and prints the end-user organization's logo onto the plastic
card. The ``card holder'' is the user whose name is printed on the
outside of the card at the same time it is written to a file inside the
card file system.\\

The following life-cycle is typical for single-application smart cards:
\begin{itemize}
  \item Control structure(behavioral) development customization.
  \item Data structure development.
  \item Plastic card manufacturing.
  \item Chip production and testing.
  \item Smart card manufacturing.
  \item Control structure(behavioral) customization.
  \item Data structure customization.
  \item Application preparation.
  \item Use.
  \item Invalidation.
\end{itemize}

The internal architecture development stage is used to develop the
control and data structures that is best suited to the requirements of
applications that will use the card. Some structures might not be
architected, and might be left open to future application developments
and extensions. Once an architected solution is found, it defines a
permanent configuration specifying the file directory structure, and
card behavior: it becomes a solution model. This solution model is a
solution for a given problem (for example, a GSM solution model or a
company solution model) that can be reproduced and manufactured.
Multiple card manufacturers can reproduce the solution model if the
solution specification follows interoperability characteristics.\\

The next requirement is for the chip for a card to be manufactured and
given a unique serial number that identifies characteristics of
manufacture (IC lot number, wafer number, die number, test equipment
numbers, etc.). The chip then undergoes a self-test process that
results either in an invalid, erased chip, or a valid chip with no
non-user modes enabled. After this stage, there is no access to card
memory addresses except under the control of the card operating system.
Tamper-resistance circuitry has been activated.\\

Independent of chip manufacture, plastic card manufacture occurs.\\

During the smart card manufacturing stage, a chip is embedded into the
plastic card. The card can be printed with the issuer's logo at this
point or later.\\

The smart card preparation stage creates the externally-downloadable
control structure and downloads it to the smart card -- changing the
basic behavior of the smart card, specifying the file structure that
will exist on the card, associating rights, keys, and privileges with
directories and files through various primitive card access and control
operations (see Security Architecture), as well as setting basic card
control parameters such as number of allowed invalid access attempts.\\

If the smart card preparation stage is not integrated with the
application preparation stage, cards need to be transported from the
card manufacturer to the end-user. In this case, the keys written at
the smart card preparation stage are termed ``transport keys'' or
``initialization keys''. These keys are sent to the card issuer by the
manufacturer separately from the cards. They prevent modification of
the cards during transport from the manufacturer to the issuer. Once in
possession of the transport keys, the issuer can access the card to
perform the application preparation stage.\\

For security purposes, the card manufcaturer can permanently set rights
and privileges before cards can move to the application preparation
stage.\\

The application preparation stage initializes the card password(s), the
data in the card file system, and optionally writes the user's name on
the card. Some cards accept dynamic update, creation and deletion of
files in the file structure after the application preparation stage
(depending on the card preparation stage), but others enforce a fixed
file structure once application preparation has taken place.\\

Then follows a stage of active card usage, during which the operations
defined in the card preparation stage may be performed using the
passwords defined during the card preparation stage. Usually, passwords
and keys may be changed during the usage stage.\\

During the final stage of the card lifecycle, the card is invalidated.
Cards should not be re-used in order to avoid any problems in
organizational policy that might arise from having two users associated
over time with the same card serial number. Invalidation of
multi-application cards can be complex if internal applications are
tied to the exterior appearance of the card: if one application is
invalidated, then the issuer logo on the card may be confusing for the
user. If the entire card must be invalidated in order to invalidate one
expired application, then the user must obtain a new card in order to
maintain access to the remaining applications.


\section{Standards}

See ``The Smart Card -- A Standardized Security Device Dedicated to
Public Cryptology'' (by Louis Claude Guillou, Michel Ugon, and
Jean-Jacques Quisquater, in Gustavus Simmons, ``Contemporary Cryptology
-- The Science of Information Integrity'', IEEE Press, pp.561-614,
1992) for an excellent discussion of the international status of smart
card standardization which is still fairly current. Several parallel,
overlapping ISO standardization efforts exist:
\begin{enumerate}
  \item ISO/IEC JTC1/SC17/WG4 (Information Systems --- Identification Cards
--- Integrated Circuit Cards with Contacts) has produced specifications of
the physical characteristics ISO 7816-1,2 and physical communication
protocols ISO 7816-3 that are International Standards. The specification of
the file system layer ISO 7816-4 is still a Committee Draft, but is
currently undergoing balloting. Other standards are still in Working Draft
stage.\\

The following standards apply for ISO SC17/WG4 contact cards:\\

Standard        Scope\\

7816-6          Common Data Elements (unpublished)\\

IS 7816-5       Registration for Applications\\

CD 7816-4.2     Commands for Interchange\\

IS 7816-3       Transmission Protocol\\

IS 7816-2       Contact Location\\

IS 7816-1       Physical Characteristics\\

As long as the lower layers conform to ISO 7816-4 and the transmission
protocol is T=0 (asynchronous, half-duplex character transmission
protocol) or T=1 (asynchronous, half-duplex block transmission
protocol), the link between the card reader and the host system can be
implemented via RS232, PCMCIA, Infra-Red, etc.

  \item ISO TC68/SC6 (Banking -- Financial Transaction Cards, Related
Media and Operations) has two working groups: WG5 (Messages Exchanged
with Integrated Circuit Cards) has produced a Draft International
Standard and a Working Draft ISO 9992 dealing with messages and data
elements. WG7 (Security Architecture of Banking Systems Using
Integrated Circuit Cards) has produced a Draft International Standard
and a Working Draft covering card life cycle, keys, and algorithms.

  \item ISO/IEC JTC1/SC17/WG8 (Information Systems --- Identification Cards
--- Contactless Integrated Circuit Cards) has not produced standards to
date in the ISO 10536 series.

  \item Specific ANSI X3B10.x groups cover contactless cards. The ANSI
X3B10.1 group covers IC cards with contacts.ANSI is circulating a
slightly modified version of ISO 7816-3.

  \item The FIPS on Security Requirements for Cryptographic Modules, which
is in draft status, specifies general physical, functional, and process
requirements for four increasing levels of security, but does not
specify interfaces or protocols. This standard applies to government
applications using cryptographic functions of smart cards.

  \item There are a few other standards bodies that have developed smart
card standards for specific industry segments or applications.
\end{enumerate}

Due to the state of standardization, it is not possible to specify a
standard interface to smart cards or smart card devices except at the
physical layer. The draft for the file system layer, however, at least
provides us with a general picture. Based on this, it is possible to
specify an abstract usage model and high-level interface that should be
implementable with minimal software ``glue'' using cards from most
vendors.


\section{Why Smart Cards In A Company?}

We have now seen the functions smart cards are capable of performing.
What value do these present for a company in particular?\\

There are three basic potential areas of value: two-factor
authentication, secure storage, and encryption and key generation.

\subsection{Two-Factor Authentication}

Without smart cards, a company uses one-factor authentication: users
authenticate by proving they know a secret that is shared with the
company Security Service, i.e., the user's password. Passwords are
known not to be very secure, as they can be lost, stolen, shared, or
guessed fairly easily.\\

By using smart cards to store each user's long-term company key, the
cards introduce a second authentication factor: users now must not only
prove they know a secret (the password used to gain access to the
card), but they must also prove they have physical possession of the
smart card (by using the password and successfully retrieving the
long-term key).\\

The use of two-factor authentication dramatically improves security.
Cards limit vulnerability to sharing, since the card can be in the
physical possession of only one individual at a time. They also
effectively prevent vulnerability to guessing, since the long-term key
stored on the card is a 56-bit random number rather than a password.
Cards can be lost or stolen, but without the accompanying card-access
password, they are not usable. Should the card be lost or stolen, the
owner is highly motivated to report the loss promptly, as the owner
will be unable to access the computer system without the card.

\subsection{Secure Storage}

The second potential value that smart cards present to a company is
their secure storage capability. This is utilized in two-factor
authentication for storing a user's long-term key. It would also be
feasible to decrypt and store company credentials on the card, hiding
them from the host system and from attackers on the host. Company
applications could also make use of secure storage on the card for
application-specific information.

\subsection{Encryption and Key Generation}

The third potential value to a company is the encryption and key
generation capabilities that smart cards have. Licensing agreements for
public key technology are typically much less restrictive when the
technology is confined to a physical device. By storing the user's
private keys in secure card storage, and using card encryption
capabilities to generate authentication information from the keys, the
encryption technology need not be implemented in the host system. The
cards' ability to generate random numbers may also be useful as a
source for keys.
